\section{数据质量评估与反馈闭环}

数据管线不是一次性脚本，而是一个需要闭环优化的系统。训练前可以用静态统计检查语料，例如语言分布、重复率、平均文档长度、域名分布、许可证比例、PII 命中率、毒性分数、代码/数学/学术文本比例；训练中可以看 loss 曲线和 domain-level loss；训练后再用 evaluation suite 反推哪些能力缺口对应哪些数据缺口。

\begin{importantbox}{从 eval 回到 data 的闭环}
\begin{enumerate}
    \item 如果代码 benchmark 弱，检查代码数据比例、license 过滤是否过严、是否缺少 issue/PR/文档数据。
    \item 如果数学推理弱，检查是否有足够高质量 LaTeX、竞赛题、步骤化解答和去重后的合成数据。
    \item 如果多语能力弱，检查 language ID 阈值是否把 code-switching 和低资源语言误删。
    \item 如果模型幻觉强，检查高权威来源比例、重复低质网页、知识密度和事实性过滤。
    \item 如果安全性差，检查 harmful data、jailbreak data、refusal data 和 policy data 的覆盖及标注质量。
\end{enumerate}
\end{importantbox}

\begin{knowledgebox}{数据审计指标}
\begin{longtable}{L{0.22\textwidth}L{0.34\textwidth}L{0.34\textwidth}}
\toprule
指标 & 看什么 & 为什么重要 \\
\midrule
\endhead
Domain mix & web、books、code、academic、forums、Q\&A、news 的比例。 & 决定模型知识和风格来源。 \\
Duplication rate & exact/near duplicate 的比例和重复 token 权重。 & 重复会浪费 compute，也可能导致 memorization。 \\
Document quality score & 规则分、classifier 分、人审样本。 & 过滤策略是否过强或过弱。 \\
License coverage & permissive、unknown、restricted、licensed 的比例。 & 决定法律和发布风险。 \\
PII/toxicity rate & 隐私和有害内容命中率。 & 决定安全和合规处理需求。 \\
Benchmark overlap & 与 eval/test 集的 exact/fuzzy overlap。 & 避免评价污染和能力高估。 \\
\bottomrule
\end{longtable}
\end{knowledgebox}

\begin{warningbox}{数据闭环也可能过拟合}
如果只根据现有 benchmark 补数据，模型会越来越适合 benchmark，而不一定更适合真实用户。数据迭代应该同时看 held-out eval、fresh eval、private eval、真实使用反馈和安全审计。
\end{warningbox}

\subsection{本章小结}

Lecture 13 的后半部分说明，同样叫“web data”可以有完全不同的数据哲学：C4 偏规则过滤，RefinedWeb/FineWeb 偏大规模 web，DCLM 偏模型质量分类器，Nemotron-CC 偏保留更多 token 和合成增强，CommonPile 偏许可可审计。数据集不是自然物，而是价值判断和工程约束的产物。

\begin{importantbox}{Data I 的根本观点}
如果说前几讲讨论的是“如何更有效地训练模型”，这一讲讨论的是“训练目标从哪里来”。数据决定模型看见什么语言、知识、代码、价值观、错误、重复和法律风险。架构可以复现，GPU 可以租，优化器可以开源；但高质量、合规、目标明确的数据管线最难复制。
\end{importantbox}

\begin{knowledgebox}{和下一讲的连接}
Data I 主要讲来源、法律边界和经典/现代数据集谱系。下一步自然是更细的数据处理问题：如何做 dedup、contamination removal、PII filtering、quality classifiers、data mixing、curriculum、post-training data 和 evaluation-driven data iteration。
\end{knowledgebox}

